\section{Discussion and applications}
\label{sec:discussion}

\paragraph{Skill-conditioned decision support.} The most direct application is
scoring that adapts to a declared skill level. Because the difficulty term is
shared and only the skill offset changes (Section~\ref{sec:model}), a single
identified model yields a family of suitability surfaces---beginner through
expert---at no additional fitting cost, each falling back invisibly to the
population curve where the latent structure is not identified. The result is a
verdict that answers ``hard for whom,'' not just ``hard on average.''

\paragraph{An intrinsic-risk primitive for underwriting.} Parametric insurance for
weather-dependent leisure needs to price the intrinsic risk of a site, not the
muddled suitability of whoever currently frequents it. The difficulty atlas is a
natural credibility primitive for such pricing: a difficulty coordinate per site
and condition, anchored to physics, reproducible from a cohort hash, and
independent of the clientele that a suitability curve would otherwise bake in. We
note the connection without building it here.

\paragraph{Limitations.} The paper's central limitation is deliberate and stated
plainly: its validation is a \emph{synthetic recovery study}, not an empirical
evaluation on real riders. We show that when data are generated from the model,
the estimator recovers the ground truth; we do \emph{not} show that real
behavioural outcomes follow the model, because no real cohort is analysed here.
Several modelling simplifications accompany this scope. The discrimination $a$ is
a single scalar rather than a function of the condition; difficulty is modelled
one physical dimension at a time rather than jointly over the full condition
vector (wind${\times}$wave${\times}$cold), which would strengthen realism at the
cost of harder identification; and estimation is point-based marginal maximum
likelihood, so the reported skills and difficulties are point estimates rather
than full posteriors. Finally, identification requires a connected
rider${\times}$condition graph (Section~\ref{sec:graph}); in real deployments many
site${\times}$condition cells will not meet that bar and will---correctly---return
the single-curve fallback rather than a latent fit.

\paragraph{Future work.} Three extensions follow directly. First, \emph{empirical
validation}: fit the model to real field cohorts as they accrue, and test the
model's structural assumptions (for instance, whether the marginal identity of
Section~\ref{sec:marginal} holds on real data, and whether the held-out skill
gain persists out of sample). Second, \emph{full Bayesian inference} with a
Gaussian-process difficulty, giving calibrated uncertainty bands on both the atlas
and the skill posteriors. Third, \emph{multi-dimensional difficulty}, jointly
identifying $\diff$ over the condition vector, which turns the per-dimension atlas
into a genuine multi-driver difficulty surface. Each is compatible with the
identification and governance framework established here.
